% Options for packages loaded elsewhere
% Options for packages loaded elsewhere
\PassOptionsToPackage{unicode}{hyperref}
\PassOptionsToPackage{hyphens}{url}
\PassOptionsToPackage{dvipsnames,svgnames,x11names}{xcolor}
%
\documentclass[
  spanish,
  letterpaper,
  conference]{IEEEtran}
\usepackage{xcolor}
\usepackage{amsmath,amssymb}
\setcounter{secnumdepth}{5}
\usepackage{iftex}
\ifPDFTeX
  \usepackage[T1]{fontenc}
  \usepackage[utf8]{inputenc}
  \usepackage{textcomp} % provide euro and other symbols
\else % if luatex or xetex
  \usepackage{unicode-math} % this also loads fontspec
  \defaultfontfeatures{Scale=MatchLowercase}
  \defaultfontfeatures[\rmfamily]{Ligatures=TeX,Scale=1}
\fi
\usepackage{lmodern}
\ifPDFTeX\else
  % xetex/luatex font selection
\fi
% Use upquote if available, for straight quotes in verbatim environments
\IfFileExists{upquote.sty}{\usepackage{upquote}}{}
\IfFileExists{microtype.sty}{% use microtype if available
  \usepackage[]{microtype}
  \UseMicrotypeSet[protrusion]{basicmath} % disable protrusion for tt fonts
}{}
\makeatletter
\@ifundefined{KOMAClassName}{% if non-KOMA class
  \IfFileExists{parskip.sty}{%
    \usepackage{parskip}
  }{% else
    \setlength{\parindent}{0pt}
    \setlength{\parskip}{6pt plus 2pt minus 1pt}}
}{% if KOMA class
  \KOMAoptions{parskip=half}}
\makeatother
% Make \paragraph and \subparagraph free-standing
\makeatletter
\ifx\paragraph\undefined\else
  \let\oldparagraph\paragraph
  \renewcommand{\paragraph}{
    \@ifstar
      \xxxParagraphStar
      \xxxParagraphNoStar
  }
  \newcommand{\xxxParagraphStar}[1]{\oldparagraph*{#1}\mbox{}}
  \newcommand{\xxxParagraphNoStar}[1]{\oldparagraph{#1}\mbox{}}
\fi
\ifx\subparagraph\undefined\else
  \let\oldsubparagraph\subparagraph
  \renewcommand{\subparagraph}{
    \@ifstar
      \xxxSubParagraphStar
      \xxxSubParagraphNoStar
  }
  \newcommand{\xxxSubParagraphStar}[1]{\oldsubparagraph*{#1}\mbox{}}
  \newcommand{\xxxSubParagraphNoStar}[1]{\oldsubparagraph{#1}\mbox{}}
\fi
\makeatother


\usepackage{longtable,booktabs,array}
\newcounter{none} % for unnumbered tables
\usepackage{calc} % for calculating minipage widths
% Correct order of tables after \paragraph or \subparagraph
\usepackage{etoolbox}
\makeatletter
\patchcmd\longtable{\par}{\if@noskipsec\mbox{}\fi\par}{}{}
\makeatother
% Allow footnotes in longtable head/foot
\IfFileExists{footnotehyper.sty}{\usepackage{footnotehyper}}{\usepackage{footnote}}
\makesavenoteenv{longtable}
\usepackage{graphicx}
\makeatletter
\newsavebox\pandoc@box
\newcommand*\pandocbounded[1]{% scales image to fit in text height/width
  \sbox\pandoc@box{#1}%
  \Gscale@div\@tempa{\textheight}{\dimexpr\ht\pandoc@box+\dp\pandoc@box\relax}%
  \Gscale@div\@tempb{\linewidth}{\wd\pandoc@box}%
  \ifdim\@tempb\p@<\@tempa\p@\let\@tempa\@tempb\fi% select the smaller of both
  \ifdim\@tempa\p@<\p@\scalebox{\@tempa}{\usebox\pandoc@box}%
  \else\usebox{\pandoc@box}%
  \fi%
}
% Set default figure placement to htbp
\def\fps@figure{htbp}
\makeatother



\ifLuaTeX
\usepackage[bidi=basic,shorthands=off]{babel}
\else
\usepackage[bidi=default,shorthands=off]{babel}
\fi
\ifLuaTeX
  \usepackage{selnolig} % disable illegal ligatures
\fi


\setlength{\emergencystretch}{3em} % prevent overfull lines

\providecommand{\tightlist}{%
  \setlength{\itemsep}{0pt}\setlength{\parskip}{0pt}}



 
\usepackage[numbers]{natbib}
\bibliographystyle{IEEEtran}


\usepackage{multirow}
\usepackage{xurl}
\usepackage{graphicx}
\AtBeginDocument{\renewcommand{\tablename}{Tabla}}
\newcommand{\notatabla}[1]{\par\vspace{2pt}\parbox{\linewidth}{\footnotesize\emph{Nota.} #1}}
\author{\IEEEauthorblockN{Fiorella Yannet Paredes Coaguila, Mario Wilfredo Ramirez Puma,\\ Jefry Erick Quispe Ramos y Paolo Fontanil Barrera}
\IEEEauthorblockA{Universidad Nacional del Altiplano -- Puno\\
Escuela Profesional de Ingeniería Estadística e Informática\\
Puno, Perú}}
\AtBeginDocument{\setlength{\parskip}{0pt}\setlength{\parindent}{1em}}
\makeatletter
\@ifpackageloaded{caption}{}{\usepackage{caption}}
\AtBeginDocument{%
\ifdefined\contentsname
  \renewcommand*\contentsname{Tabla de contenidos}
\else
  \newcommand\contentsname{Tabla de contenidos}
\fi
\ifdefined\listfigurename
  \renewcommand*\listfigurename{Listado de Figuras}
\else
  \newcommand\listfigurename{Listado de Figuras}
\fi
\ifdefined\listtablename
  \renewcommand*\listtablename{Listado de Tablas}
\else
  \newcommand\listtablename{Listado de Tablas}
\fi
\ifdefined\figurename
  \renewcommand*\figurename{Figura}
\else
  \newcommand\figurename{Figura}
\fi
\ifdefined\tablename
  \renewcommand*\tablename{Tabla}
\else
  \newcommand\tablename{Tabla}
\fi
}
\@ifpackageloaded{float}{}{\usepackage{float}}
\floatstyle{ruled}
\@ifundefined{c@chapter}{\newfloat{codelisting}{h}{lop}}{\newfloat{codelisting}{h}{lop}[chapter]}
\floatname{codelisting}{Listado}
\newcommand*\listoflistings{\listof{codelisting}{Listado de Listados}}
\makeatother
\makeatletter
\makeatother
\makeatletter
\@ifpackageloaded{caption}{}{\usepackage{caption}}
\@ifpackageloaded{subcaption}{}{\usepackage{subcaption}}
\makeatother
\usepackage{bookmark}
\IfFileExists{xurl.sty}{\usepackage{xurl}}{} % add URL line breaks if available
\urlstyle{same}
% fallback for those not using the hyperref driver hyperxmp:
\makeatletter
\@ifundefined{xmpquote}{\newcommand{\xmpquote}[1]{#1}}{}
\makeatother
\hypersetup{
  pdftitle={Desigualdades en la atención prenatal y el parto institucional entre mujeres quechuahablantes{,} aimarahablantes y castellanohablantes del departamento de Puno{,} Perú: análisis de la ENDES 2014--2025},
  pdflang={es},
  colorlinks=true,
  linkcolor={blue},
  filecolor={Maroon},
  citecolor={Blue},
  urlcolor={Blue},
  pdfcreator={LaTeX via pandoc}}


\title{Desigualdades en la atención prenatal y el parto institucional
entre mujeres quechuahablantes, aimarahablantes y castellanohablantes
del departamento de Puno, Perú: análisis de la ENDES 2014--2025}
\author{}
\date{}
\begin{document}
\maketitle


\begin{abstract}
\textit{Objetivo.} Describir las diferencias en la atención prenatal y el parto institucional según lengua materna entre mujeres de Puno y estimar asociaciones ajustadas por características sociodemográficas con la ENDES 2014--2025. \textit{Métodos.} Estudio observacional de cortes transversales repetidos. Se incluyeron 5,684 mujeres con un nacimiento en los 60 meses previos a la encuesta, considerando el más reciente. Los desenlaces fueron seis o más controles prenatales, inicio del control en el primer trimestre y parto institucional. Se estimaron porcentajes ponderados y modelos logísticos para muestras complejas ajustados por año, edad, residencia, educación, riqueza y número de hijos; se calcularon probabilidades ajustadas, diferencias absolutas y la interacción lengua $\times$ periodo. El análisis completo se reproduce con R y Quarto. \textit{Resultados.} Participaron 1,739 castellanohablantes, 2,639 quechuahablantes y 1,306 aimarahablantes. Frente a castellanohablantes, la razón de momios ajustada (ORa) de parto institucional fue 0,59 (0,43--0,82) en quechuahablantes y 0,66 (0,45--0,96) en aimarahablantes, equivalentes a diferencias ajustadas de $-$3,7 y $-$2,9 puntos porcentuales; el ajuste explicó cerca de la mitad de la asociación cruda. La mayor diferencia absoluta correspondió al inicio temprano entre aimarahablantes ($-$8,8 puntos; ORa = 0,68). Las asociaciones no disminuyeron significativamente entre 2014--2019 y 2020--2025 y fueron robustas a una definición alternativa de parto institucional y al tratamiento de la varianza. \textit{Conclusiones.} Persisten desventajas asociadas a la lengua materna después del ajuste sociodemográfico, sobre todo en parto institucional e inicio oportuno del control prenatal.
\end{abstract}
\begin{IEEEkeywords}
salud materna; control prenatal; parto institucional; lenguas indígenas; desigualdad en salud; ENDES; Puno; reproducibilidad.
\end{IEEEkeywords}

\section{Introducción}\label{introducciuxf3n}

La atención materna oportuna, continua y culturalmente pertinente
constituye un componente fundamental de los sistemas de salud. El
seguimiento durante la gestación permite identificar señales de alarma,
promover medidas preventivas y establecer vínculos entre las mujeres y
los servicios sanitarios. La Organización Mundial de la Salud recomienda
un modelo centrado en la experiencia positiva del embarazo y en al menos
ocho contactos de atención prenatal \citep{oms2016}. La cobertura, por
sí sola, no equivale a calidad: también importan la oportunidad, la
información, el trato y la capacidad de respuesta del establecimiento.

Las desigualdades de acceso se expresan de forma distinta según
condiciones sociales y culturales. En América Latina, estudios
comparativos han identificado que las poblaciones indígenas enfrentan
desventajas en diversas intervenciones de salud reproductiva y materna,
especialmente aquellas que dependen de acudir a establecimientos y
disponer de recursos para recibir atención \citep{mesenburg2018}. La
lengua hablada constituye una aproximación útil ---aunque incompleta---
a experiencias lingüísticas y sociales relacionadas con la atención
sanitaria, y ha sido empleada para documentar inequidades maternas en
Bolivia, Guatemala, México y Perú \citep{armenta2019}. Su uso exige
distinguir la variable observada de conceptos más amplios como la
pertenencia étnica o la autoidentificación \citep{armenta2020}.

Para Perú, se ha documentado que las diferencias socioeconómicas
persisten cuando se examina la cantidad y calidad de controles
prenatales \citep{enriquez2022}. Investigaciones centradas en el parto
institucional también encontraron asociaciones con área de residencia,
nivel educativo, riqueza y aseguramiento \citep{rodrigo2023, ramos2023}.
Tales antecedentes orientan la elección de covariables y desaconsejan
interpretar automáticamente las diferencias descriptivas entre lenguas
como efectos causales del idioma.

En el departamento de Puno, la coexistencia de castellano, quechua y
aimara hace pertinente un análisis específico. El estudio cualitativo de
Samuel \citep{samuel2016} muestra que la implementación de políticas
para la atención intercultural ha requerido participación social y
vigilancia de la calidad de los servicios. La experiencia de adaptación
cultural de la atención del parto en Ayacucho \citep{gabrysch2009}
permite entender que los cambios organizacionales, la comunicación y las
prácticas respetuosas pueden contribuir a reducir barreras percibidas.
Los trabajos etnográficos y cualitativos en la Amazonía peruana
\citep{delmastro2021, mejia2025} ofrecen contextos valiosos, pero no
deben confundirse con evidencia causal ni extrapolarse de forma
automática a Puno.

Aunque se dispone de informes anuales de la Encuesta Demográfica y de
Salud Familiar (ENDES) \citep{inei2024, inei2025}, comparar de manera
sistemática a mujeres según lengua materna en Puno durante 2014--2025
permite examinar patrones de largo plazo y asociaciones ajustadas. El
valor del estudio radica en generar evidencia local, totalmente
reproducible, para orientar el seguimiento de desigualdades lingüísticas
y socioeconómicas en servicios maternos.

\section{Objetivos}\label{objetivos}

\emph{Objetivo general.} Analizar las diferencias en la cobertura de
atención prenatal y parto institucional entre mujeres
castellanohablantes, quechuahablantes y aimarahablantes de Puno
utilizando datos de la ENDES 2014--2025 y estimar su asociación con la
lengua materna después de ajustar por características sociodemográficas.

\emph{Objetivos específicos.} (a) Describir la composición
sociodemográfica de los grupos lingüísticos; (b) estimar porcentajes
anuales ponderados de seis o más controles prenatales, inicio temprano
del control prenatal y parto institucional; (c) estimar las asociaciones
entre lengua materna y los tres indicadores, en escala relativa y
absoluta; (d) evaluar si dichas asociaciones cambiaron entre 2014--2019
y 2020--2025; y (e) documentar el análisis de forma reproducible con R y
Quarto.

\section{Materiales y métodos}\label{materiales-y-muxe9todos}

\subsection{Diseño y fuente de
información}\label{diseuxf1o-y-fuente-de-informaciuxf3n}

Estudio cuantitativo, observacional y analítico de cortes transversales
repetidos, basado en microdatos de acceso público de doce ediciones de
la ENDES entre 2014 y 2025 \citep{inei2024, inei2025}. La información
fue armonizada por año, identificando registros de mujeres, antecedentes
obstétricos del nacimiento más reciente y variables del diseño muestral.
Se integró una base multianual con identificadores de conglomerado y
estrato diferenciados por año. Los resultados describen asociaciones
para los dominios investigados, no seguimientos longitudinales de las
mismas mujeres.

\subsection{Población y muestra}\label{poblaciuxf3n-y-muestra}

De 12~620 mujeres de 15 a 49 años entrevistadas en Puno, 5~689
informaron al menos un nacimiento en los 60 meses previos a la
entrevista (6~591 nacimientos). Para evitar contar varios partos de la
misma mujer en una encuesta, se utilizó el nacimiento más reciente. Se
excluyeron 5 mujeres cuya lengua materna no era castellano, quechua ni
aimara o carecía de dato, de modo que la muestra analítica comprendió
5~684 mujeres: castellano (1~739), quechua (2~639) y aimara (1~306)
(Fig.~\ref{fig-flujo}).

\begin{figure}[!t]

\centering{

\includegraphics[width=1\linewidth,height=\textheight,keepaspectratio]{informe_files/figure-pdf/fig-flujo-1.pdf}

}

\caption{\label{fig-flujo}Diagrama de flujo de selección de la muestra
analítica, Puno, ENDES 2014--2025.}

\end{figure}%

\subsection{Variables de estudio}\label{variables-de-estudio}

La exposición principal fue la lengua materna, reconstruida a partir de
la variable S119. Se armonizaron los códigos entre ediciones: en
2014--2016, castellano = 1, quechua = 2 y aimara = 3; en 2017--2025,
quechua = 1, aimara = 2 y castellano = 10. Esta recodificación es
indispensable para evitar clasificación errónea en las comparaciones
temporales.

Los desenlaces binarios fueron: (1) seis o más controles prenatales, con
al menos seis visitas registradas en M14; (2) control temprano, definido
como primera visita durante los meses 1--3 según M13, clasificando la
ausencia de controles como no temprana; y (3) parto institucional,
construido con el lugar del parto (M15) y la atención por personal
profesional (M3A--M3C). Se clasificaron como no institucionales los
partos en domicilio propio o de partera; las categorías no clasificables
(M15 = 96) se trataron como faltantes. Dado que lugar del parto y
atención calificada no son necesariamente un mismo indicador, se repitió
el análisis con una definición basada únicamente en el lugar del parto.
Como indicador complementario, alineado con la recomendación vigente de
la OMS \citep{oms2016}, se analizó además la proporción con ocho o más
controles.

Como variables de ajuste se incluyeron edad en años (V012), residencia
urbana/rural (V025), nivel educativo (V106), quintil de riqueza (V190),
número de hijos nacidos vivos (V201) y año de encuesta. Educación,
riqueza y residencia se incorporaron como categorías; edad y número de
hijos, como términos numéricos.

\subsection{Diseño muestral y análisis
estadístico}\label{diseuxf1o-muestral-y-anuxe1lisis-estaduxedstico}

Los análisis utilizaron factores de ponderación V005/1 000 000,
conglomerados V001 y estratos basados en V022 junto con el departamento
V024. Los códigos de conglomerados y estratos se distinguieron entre
años. El diseño se construyó con los registros nacionales de cada
edición (434~993 mujeres) y luego se seleccionó el dominio de estudio,
evitando eliminar unidades del diseño de forma prematura. En el análisis
agrupado, los pesos se dividieron entre 12, de modo que las estimaciones
combinadas representan un año promedio del periodo; este reescalamiento
no modifica los coeficientes.

Se estimaron porcentajes por año y lengua con intervalos de confianza
del 95 \% por el método logit, con un diseño por encuesta. Para cada
desenlace se ajustaron modelos logísticos binarios con \texttt{svyglm()}
y familia \texttt{quasibinomial()} del paquete \texttt{survey} de R
\citep{lumley2004, rcore}, tomando como referencia a las
castellanohablantes. Se informan razones de momios crudas (ORc, con la
lengua materna como única variable) y ajustadas (ORa), intervalos del 95
\% y valores \(p\); la comparación entre ambas indica qué parte de la
asociación se explica por las covariables. Como los desenlaces fueron
frecuentes y los OR exageran la magnitud relativa en esa situación, se
calcularon también probabilidades ajustadas por estandarización marginal
(asignando a cada mujer cada lengua y promediando con los pesos del
diseño) y sus diferencias en puntos porcentuales, con errores estándar
por el método delta.

Para examinar si las desigualdades cambiaron en el tiempo, se añadió a
cada modelo la interacción lengua \(\times\) periodo (2014--2019 frente
a 2020--2025), evaluada con una prueba de Wald ajustada al diseño. En
algunos años se detectaron estratos con una sola unidad primaria de
muestreo; el análisis principal utilizó la opción
\texttt{survey.lonely.psu\ =\ "adjust"} y, como análisis de
sensibilidad, se repitió con \texttt{"average"}. La significancia se
evaluó al 5 \%, sin interpretar la ausencia de significancia como
ausencia de diferencia. El informe sigue las recomendaciones STROBE
\citep{vonelm2007, vandenbroucke2007}.

\subsection{Reproducibilidad}\label{reproducibilidad}

El informe se escribió en Quarto \citep{quarto} y se ejecutó en Positron
con R. Cada cifra, tabla y figura de este documento se calcula al
renderizar el archivo \texttt{informe.qmd} a partir de dos archivos
livianos: la base analítica de 5~684 mujeres y el marco de diseño
muestral nacional; ningún resultado se transcribió a mano. La
preparación de esos archivos desde los microdatos originales del INEI se
documenta en los scripts numerados 01--51 del proyecto. La sección final
detalla el repositorio y las versiones del software.

\subsection{Consideraciones éticas}\label{consideraciones-uxe9ticas}

El estudio utilizó microdatos anonimizados de acceso público difundidos
por el Instituto Nacional de Estadística e Informática (INEI), sin
contacto ni intervención sobre las participantes. Los resultados se
presentan únicamente de forma agregada.

\section{Resultados}\label{resultados}

\subsection{Composición de la muestra y características
sociodemográficas}\label{composiciuxf3n-de-la-muestra-y-caracteruxedsticas-sociodemogruxe1ficas}

De las 5~684 mujeres elegibles, el 46,4 \% fueron quechuahablantes, el
30,6 \% castellanohablantes y el 23,0 \% aimarahablantes; estos
porcentajes se calcularon sobre recuentos sin ponderar y no representan
la composición poblacional. La Tabla \ref{tab:socio} y la
Fig.~\ref{fig-socio} resumen las características ponderadas.

\begin{table}[!t]
\centering
\caption{Características sociodemográficas por lengua materna}
\label{tab:socio}
\small
\setlength{\tabcolsep}{4pt}
\begin{tabular*}{\columnwidth}{@{\extracolsep{\fill}}lrrr@{}}
\toprule
Característica & Castellano & Quechua & Aimara \\
\midrule
$n$ (sin ponderar) & 1,739 & 2,639 & 1,306 \\
Edad media (años) & 28,5 & 30,8 & 31,0 \\
Residencia rural (\%) & 24,5 & 49,5 & 56,7 \\
Educación superior (\%) & 43,4 & 17,8 & 16,8 \\
Dos quintiles inferiores (\%) & 53,0 & 79,5 & 79,7 \\
Hijos nacidos vivos, media & 1,88 & 2,58 & 2,56 \\
\bottomrule
\end{tabular*}
\notatabla{Estimaciones ponderadas según el diseño muestral de la ENDES.}
\end{table}

\begin{figure}[!t]

\centering{

\includegraphics[width=1\linewidth,height=\textheight,keepaspectratio]{informe_files/figure-pdf/fig-socio-1.pdf}

}

\caption{\label{fig-socio}Características sociodemográficas
seleccionadas por lengua materna (porcentajes ponderados), Puno, ENDES
2014--2025.}

\end{figure}%

Las mayores diferencias se concentraron en residencia, escolaridad y
riqueza. La residencia rural alcanzó 49,5 \% entre quechuahablantes y
56,7 \% entre aimarahablantes, frente a 24,5 \% en castellanohablantes.
La educación superior mostró el patrón inverso: 43,4 \% en castellano,
17,8 \% en quechua y 16,8 \% en aimara. Esta heterogeneidad justifica
presentar, además de las diferencias descriptivas, comparaciones
ajustadas.

\subsection{Evolución anual de los
indicadores}\label{evoluciuxf3n-anual-de-los-indicadores}

La Tabla \ref{tab:anual} y la Fig.~\ref{fig-tendencias} presentan los
porcentajes ponderados para los doce años disponibles. La cobertura de
seis o más visitas y el parto institucional alcanzan niveles elevados en
los últimos años, mientras que el control temprano mantiene diferencias
más visibles. Los intervalos de confianza anuales son amplios, por lo
que las fluctuaciones de un año a otro no deben interpretarse como
cambios significativos.

\begin{table*}[!t]
\centering
\caption{Indicadores de atención materna por año y lengua materna (\%)}
\label{tab:anual}
\small
\setlength{\tabcolsep}{6pt}
\begin{tabular}{lrrrrrrrrr}
\toprule
& \multicolumn{3}{c}{$\geq$6 controles prenatales} & \multicolumn{3}{c}{Inicio en el primer trimestre} & \multicolumn{3}{c}{Parto institucional} \\
\cmidrule(lr){2-4}\cmidrule(lr){5-7}\cmidrule(lr){8-10}
Año & Cast. & Quechua & Aimara & Cast. & Quechua & Aimara & Cast. & Quechua & Aimara \\
\midrule
2014 & 82,4 & 79,4 & 68,2 & 60,5 & 59,5 & 47,2 & 91,8 & 69,3 & 77,3 \\
2015 & 81,1 & 81,2 & 75,5 & 69,7 & 56,6 & 55,5 & 84,8 & 79,4 & 77,5 \\
2016 & 79,9 & 81,9 & 78,9 & 63,9 & 62,1 & 61,0 & 93,8 & 84,6 & 85,6 \\
2017 & 82,7 & 82,0 & 75,2 & 61,5 & 65,5 & 57,4 & 97,3 & 86,2 & 91,0 \\
2018 & 80,8 & 79,0 & 75,7 & 64,1 & 57,6 & 52,8 & 96,9 & 88,8 & 84,7 \\
2019 & 87,2 & 81,2 & 83,7 & 71,5 & 62,3 & 57,5 & 95,1 & 89,0 & 90,8 \\
2020 & 90,6 & 82,3 & 88,0 & 81,2 & 66,4 & 58,4 & 95,6 & 89,9 & 94,5 \\
2021 & 82,5 & 83,1 & 78,4 & 70,9 & 60,1 & 59,2 & 98,9 & 93,7 & 88,0 \\
2022 & 86,7 & 81,4 & 87,7 & 66,8 & 62,9 & 54,0 & 99,0 & 94,6 & 95,1 \\
2023 & 88,2 & 86,5 & 90,8 & 77,6 & 66,4 & 68,2 & 100,0 & 97,5 & 99,3 \\
2024 & 90,8 & 87,2 & 78,2 & 74,1 & 68,1 & 60,8 & 97,0 & 97,4 & 93,4 \\
2025 & 87,1 & 87,1 & 89,6 & 80,8 & 67,7 & 64,3 & 98,7 & 98,4 & 96,7 \\
\bottomrule
\end{tabular}
\notatabla{Porcentajes ponderados, redondeados a un decimal, estimados con el diseño muestral de cada encuesta. Los intervalos de confianza se muestran en la Fig.~\ref{fig-tendencias}. Un valor de 100\,\% indica que todas las observaciones válidas del grupo fueron positivas y no implica certeza poblacional.}
\end{table*}

\begin{figure}[!t]

\centering{

\includegraphics[width=1\linewidth,height=\textheight,keepaspectratio]{informe_files/figure-pdf/fig-tendencias-1.pdf}

}

\caption{\label{fig-tendencias}Evolución anual de los indicadores por
lengua materna, Puno, 2014--2025: (a) seis o más controles prenatales;
(b) inicio del control en el primer trimestre; (c) parto institucional.
Líneas: porcentajes ponderados; bandas: IC del 95 \% (método logit).
Cuando la estimación es 100 \% no se muestra banda.}

\end{figure}%

En 2014, el parto institucional fue estimado en 91,8 \% para castellano,
69,3 \% para quechua y 77,3 \% para aimara, mientras que en 2025 alcanzó
98,7 \%, 98,4 \% y 96,7 \%, respectivamente. Estas diferencias entre
extremos de la serie son descriptivas y no permiten concluir causalmente
que un grupo se benefició más que otro.

Para inicio temprano, las estimaciones de 2025 fueron 80,8 \% en
castellano, 67,7 \% en quechua y 64,3 \% en aimara, con diferencias
descriptivas frente a castellano de 13,1 y 16,5 puntos porcentuales,
respectivamente.

\subsection{Modelos logísticos}\label{modelos-loguxedsticos}

La Tabla \ref{tab:ora} y la Fig.~\ref{fig-forest} muestran las
asociaciones entre lengua materna y cada desenlace, sin ajuste y
controlando por edad, residencia, educación, riqueza, número de hijos
nacidos vivos y año de encuesta. Los modelos trabajaron con \(n\) =
5~681 para seis o más controles, \(n\) = 5~683 para inicio temprano y
\(n\) = 5~645 para parto institucional.

\begin{table}[!t]
\centering
\caption{Razones de momios crudas y ajustadas por lengua materna}
\label{tab:ora}
\footnotesize
\setlength{\tabcolsep}{2pt}
\begin{tabular*}{\columnwidth}{@{\extracolsep{\fill}}llccc@{}}
\toprule
Indicador & Lengua & ORc (IC 95\,\%) & ORa (IC 95\,\%) & $p$ \\
\midrule
$\geq$6 contr. & Quechua & 0,87 (0,72--1,05) & 0,95 (0,78--1,16) & 0,613 \\
 & Aimara & 0,75 (0,59--0,94) & 0,76 (0,60--0,97) & 0,027 \\
\addlinespace
Inicio temp. & Quechua & 0,74 (0,63--0,86) & 0,85 (0,72--1,00) & 0,053 \\
 & Aimara & 0,61 (0,51--0,73) & 0,68 (0,56--0,81) & $<$0,001 \\
\addlinespace
Parto inst. & Quechua & 0,35 (0,26--0,49) & 0,59 (0,43--0,82) & 0,002 \\
 & Aimara & 0,38 (0,26--0,55) & 0,66 (0,45--0,96) & 0,029 \\
\bottomrule
\end{tabular*}
\notatabla{ORc: razón de momios cruda; ORa: ajustada por edad, residencia, educación, quintil de riqueza, número de hijos nacidos vivos y año de encuesta. Referencia: castellano. Ambos modelos usan los mismos casos y el diseño muestral; $p$ corresponde al ORa. Contr.: controles; temp.: temprano; inst.: institucional.}
\end{table}

\begin{figure}[!t]

\centering{

\includegraphics[width=1\linewidth,height=\textheight,keepaspectratio]{informe_files/figure-pdf/fig-forest-1.pdf}

}

\caption{\label{fig-forest}Razones de momios ajustadas (ORa) e IC del 95
\% de la lengua materna para cada indicador (referencia: castellano).}

\end{figure}%

Frente a las castellanohablantes, las aimarahablantes presentaron
menores momios ajustados de haber tenido seis o más controles (ORa =
0,763; IC 95 \%: 0,601--0,970; \(p\) = 0,027). Para quechua, la
estimación fue cercana a la unidad (ORa = 0,950; \(p\) = 0,613). En el
inicio temprano, el ORa fue 0,675 en aimara (\(p\) \textless{} 0,001) y
0,850 en quechua (\(p\) = 0,053), este último con evidencia limitada en
el umbral convencional del 5 \%. En el parto institucional, los dos
grupos de lengua indígena presentaron menores momios ajustados: ORa =
0,592 para quechua (IC 95 \%: 0,427--0,821) y ORa = 0,656 para aimara
(IC 95 \%: 0,449--0,957).

El ajuste por características sociodemográficas atenuó de forma
importante la asociación con el parto institucional: la razón de momios
pasó de 0,35 (cruda) a 0,59 (ajustada) en quechuahablantes y de 0,38 a
0,66 en aimarahablantes, es decir, el coeficiente se redujo en 50 \% y
57 \% en escala logarítmica. En cambio, para el inicio temprano en
aimarahablantes la atenuación fue de solo 20 \% (ORc = 0,61; ORa =
0,68), de modo que la mayor parte de esa desventaja no se explicó por
las covariables consideradas.

\subsection{Probabilidades ajustadas y diferencias
absolutas}\label{probabilidades-ajustadas-y-diferencias-absolutas}

Debido a la alta frecuencia de los desenlaces, la magnitud de los OR no
se traduce directamente en diferencias de probabilidad (Tabla
\ref{tab:abs}). En parto institucional, las diferencias ajustadas frente
a castellano fueron de \(-3{,}7\) puntos porcentuales en quechua (IC 95
\%: \(-5{,}9\) a \(-1{,}6\)) y \(-2{,}9\) en aimara. La mayor diferencia
absoluta correspondió al inicio temprano en aimarahablantes: 58,5 \%
frente a 67,2 \% en castellanohablantes, una diferencia de \(-8{,}8\)
puntos (\(-12{,}8\) a \(-4{,}7\)).

\begin{table*}[!t]
\centering
\caption{Probabilidades ajustadas y diferencias absolutas frente a castellanohablantes}
\label{tab:abs}
\small
\setlength{\tabcolsep}{5pt}
\begin{tabular}{lccccc}
\toprule
& \multicolumn{3}{c}{Probabilidad ajustada, \% (IC 95\,\%)} & \multicolumn{2}{c}{Diferencia frente a castellano, pp (IC 95\,\%)} \\
\cmidrule(lr){2-4}\cmidrule(lr){5-6}
Indicador & Castellano & Quechua & Aimara & Quechua & Aimara \\
\midrule
$\geq$6 controles & 83,7 (81,6--85,8) & 83,0 (81,4--84,7) & 79,8 (77,1--82,5) & $-$0,7 ($-$3,4 a 2,0) & $-$3,9 ($-$7,4 a $-$0,4) \\
Inicio temprano & 67,2 (64,6--69,9) & 63,7 (61,5--65,9) & 58,5 (55,4--61,5) & $-$3,5 ($-$7,1 a 0,0) & $-$8,8 ($-$12,8 a $-$4,7) \\
Parto institucional & 93,1 (91,4--94,9) & 89,4 (87,9--90,9) & 90,2 (88,4--92,0) & $-$3,7 ($-$5,9 a $-$1,6) & $-$2,9 ($-$5,5 a $-$0,4) \\
\bottomrule
\end{tabular}
\notatabla{Probabilidades obtenidas por estandarización marginal a partir de los modelos ajustados de la Tabla~\ref{tab:ora}, con los pesos del diseño. pp: puntos porcentuales.}
\end{table*}

\subsection{Cambio de las desigualdades en el tiempo y análisis de
sensibilidad}\label{cambio-de-las-desigualdades-en-el-tiempo-y-anuxe1lisis-de-sensibilidad}

La Tabla \ref{tab:periodo} muestra las ORa por periodo. En ningún
indicador la interacción lengua \(\times\) periodo alcanzó significancia
estadística (\(p\) = 0,093 para seis o más controles, 0,072 para inicio
temprano y 0,604 para parto institucional). En parto institucional, pese
a la convergencia de los porcentajes, las ORa de 2020--2025 se
mantuvieron por debajo de la unidad (0,51 en quechua y 0,47 en aimara).
En inicio temprano, la asociación en quechuahablantes fue nula en
2014--2019 (ORa = 0,99) y se hizo evidente en 2020--2025 (ORa = 0,69; IC
95 \%: 0,56--0,85).

\begin{table}[!t]
\centering
\caption{Razones de momios ajustadas por periodo}
\label{tab:periodo}
\footnotesize
\setlength{\tabcolsep}{2.2pt}
\begin{tabular*}{\columnwidth}{@{\extracolsep{\fill}}llccc@{}}
\toprule
Indicador & Lengua & 2014--2019 & 2020--2025 & $p_{\mathrm{int}}$ \\
\midrule
$\geq$6 contr. & Quechua & 1,05 (0,82--1,35) & 0,81 (0,60--1,10) & \multirow{2}{*}{0,093} \\
 & Aimara & 0,72 (0,53--0,98) & 0,85 (0,59--1,24) &  \\
\addlinespace
Inicio temp. & Quechua & 0,99 (0,78--1,24) & 0,69 (0,56--0,85) & \multirow{2}{*}{0,072} \\
 & Aimara & 0,75 (0,58--0,97) & 0,58 (0,46--0,74) &  \\
\addlinespace
Parto inst. & Quechua & 0,61 (0,42--0,89) & 0,51 (0,26--0,997) & \multirow{2}{*}{0,604} \\
 & Aimara & 0,71 (0,46--1,10) & 0,47 (0,22--0,996) &  \\
\bottomrule
\end{tabular*}
\notatabla{ORa (IC 95\,\%) frente a castellano en cada periodo, a partir de modelos con interacción lengua $\times$ periodo y las mismas covariables. $p_{\mathrm{int}}$: prueba de Wald de la interacción (2 grados de libertad).}
\end{table}

Los resultados fueron robustos (Tabla \ref{tab:sens}). Con la definición
de parto institucional basada solo en el lugar del parto, las ORa fueron
0,60 en quechua y 0,68 en aimara, prácticamente iguales a las del
análisis principal. Con el tratamiento \texttt{average} para estratos
con una sola unidad primaria, los coeficientes no cambiaron y los
intervalos se ampliaron levemente, sin modificar la significancia. Por
último, la cobertura de ocho o más controles fue considerablemente menor
que la de seis o más (64,5 \% en castellano, 60,6 \% en quechua y 61,1
\% en aimara), sin diferencias significativas por lengua.

\begin{table}[!t]
\centering
\caption{Análisis de sensibilidad e indicador complementario}
\label{tab:sens}
\small
\setlength{\tabcolsep}{3pt}
\begin{tabular*}{\columnwidth}{@{\extracolsep{\fill}}llcc@{}}
\toprule
Análisis & Lengua & ORa (IC 95\,\%) & $p$ \\
\midrule
\multicolumn{4}{l}{\textit{Parto institucional solo por lugar del parto}} \\
& Quechua & 0,60 (0,43--0,83) & 0,002 \\
& Aimara & 0,68 (0,46--0,997) & 0,048 \\
\addlinespace
\multicolumn{4}{l}{\textit{Parto institucional, varianza con \texttt{average}}} \\
& Quechua & 0,59 (0,43--0,82) & 0,002 \\
& Aimara & 0,66 (0,45--0,96) & 0,029 \\
\addlinespace
\multicolumn{4}{l}{\textit{Inicio temprano, varianza con \texttt{average}}} \\
& Quechua & 0,85 (0,72--1,00) & 0,053 \\
& Aimara & 0,68 (0,56--0,81) & $<$0,001 \\
\addlinespace
\multicolumn{4}{l}{\textit{$\geq$6 controles, varianza con \texttt{average}}} \\
& Quechua & 0,95 (0,78--1,16) & 0,613 \\
& Aimara & 0,76 (0,60--0,97) & 0,027 \\
\addlinespace
\multicolumn{4}{l}{\textit{$\geq$8 controles (recomendación OMS 2016)}} \\
& Quechua & 0,87 (0,74--1,03) & 0,104 \\
& Aimara & 0,86 (0,71--1,04) & 0,125 \\
\bottomrule
\end{tabular*}
\notatabla{Referencia: castellano. Mismas covariables que en la Tabla~\ref{tab:ora}.}
\end{table}

\section{Discusión}\label{discusiuxf3n}

El presente análisis integra doce ediciones de la ENDES para explorar
desigualdades vinculadas con la lengua materna en Puno. Su principal
contribución descriptiva consiste en presentar, en una misma serie, tres
dimensiones de atención materna comparadas entre castellano, quechua y
aimara. La contribución analítica reside en mostrar que algunas
asociaciones con la lengua materna persisten cuando se consideran
diferencias de edad, residencia, educación, riqueza, paridad y año,
particularmente para parto institucional y, entre aimarahablantes, para
el inicio temprano del control prenatal.

Las diferencias en composición sociodemográfica sugieren que la
desigualdad lingüística se superpone con desventajas territoriales y
materiales. Las mujeres aimarahablantes y quechuahablantes estaban más
representadas en áreas rurales y en los quintiles inferiores de riqueza,
patrón coherente con la literatura latinoamericana sobre inequidades
según riqueza, territorio y pertenencia étnica en salud maternoinfantil
\citep{mesenburg2018, colome2021}. La comparación entre las razones
crudas y ajustadas lo confirma: cerca de la mitad de la desventaja en
parto institucional se explicó por residencia, educación, riqueza y
paridad, pero la otra mitad persistió. En el inicio temprano entre
aimarahablantes, en cambio, el ajuste apenas modificó la asociación, lo
que apunta a factores no captados por las variables socioeconómicas,
como la distancia, la comunicación con el personal o la experiencia
previa con los servicios. No obstante, estos resultados no permiten
distinguir con precisión las rutas causales de esas desigualdades.

La asociación más consistente en ambos grupos indígenas correspondió al
parto institucional. Ello es relevante porque el acceso a servicios para
el parto depende de un entramado de distancia, transporte, costo
indirecto, experiencias de discriminación, pertinencia lingüística y
confianza en el establecimiento. La investigación de Samuel
\citep{samuel2016} en Puno destaca el papel de la vigilancia ciudadana y
la adecuación intercultural; la experiencia de Gabrysch et al.
\citep{gabrysch2009} en Ayacucho documenta la adaptación de servicios a
prácticas y preferencias locales, y estudios en la Amazonía describen
preferencias por el parto domiciliario y barreras de acceso en
comunidades indígenas \citep{delmastro2021, mejia2025}. Aunque los
contextos sean diferentes, estos estudios fundamentan la necesidad de
interpretar las cifras de cobertura junto con la calidad y el trato
percibido.

Las escalas relativa y absoluta ofrecen lecturas complementarias. Un ORa
de 0,59 corresponde a una diferencia ajustada de apenas 3--4 puntos
porcentuales cuando el parto institucional supera el 90 \%. Que la
diferencia relativa no haya disminuido entre 2014--2019 y 2020--2025,
aun cuando los porcentajes convergieron, indica que los partos que
todavía ocurren fuera de establecimientos se concentran en mujeres de
lengua indígena; este grupo residual puede requerir estrategias
focalizadas más que una expansión general de la cobertura.

En la atención prenatal, el resultado más reciente indica porcentajes
elevados de seis o más consultas en los tres grupos, mientras que el
inicio en el primer trimestre muestra las mayores diferencias absolutas,
de casi nueve puntos entre aimarahablantes. Un inicio tardío restringe
las oportunidades de orientación preventiva y de seguimiento adecuado.
La aparición de una desventaja en quechuahablantes durante 2020--2025
coincide con el periodo de interrupción de servicios asociado a la
pandemia de COVID-19, aunque este estudio no permite atribuirla a esa
causa y la prueba de interacción no fue concluyente. Asimismo, la
cobertura de ocho o más controles, estándar actual de la OMS
\citep{oms2016}, fue unos 20 puntos inferior a la de seis o más, lo que
indica que el indicador histórico sobrestima la adecuación del control
prenatal en todos los grupos.

La literatura peruana sobre cobertura y calidad prenatal
\citep{enriquez2022} respalda examinar no solamente cuántas consultas
ocurrieron, sino también su oportunidad y contenido. Asimismo, los
estudios sobre información recibida durante la atención y sobre
aseguramiento indican que los factores institucionales y
comunicacionales influyen en la utilización de servicios
\citep{quispe2023, ramos2023, perezlu2018}. Una implicancia práctica de
este trabajo es recomendar la evaluación de la oferta lingüística y de
la experiencia de las usuarias, además del monitoreo cuantitativo de la
cobertura.

Una lectura prudente exige recordar que la lengua materna no equivale
por sí misma a identidad étnica, discriminación experimentada o
competencia comunicativa con el personal de salud. La medición de
inequidades étnicas depende del criterio de clasificación empleado
\citep{armenta2020}; por ello, las asociaciones encontradas podrían
reflejar factores culturales, sociales, geográficos e institucionales no
observados. Tampoco puede asumirse que las diferencias entre 2014 y 2025
se deban a cambios en políticas específicas, dado el carácter repetido y
transversal del diseño.

\subsection{Fortalezas y limitaciones}\label{fortalezas-y-limitaciones}

Entre las fortalezas destacan la integración temporal de doce encuestas,
la armonización explícita de la codificación lingüística, el uso del
diseño muestral completo, el análisis de tres resultados de atención
materna en escalas relativa y absoluta, la verificación de la robustez
frente a definiciones alternativas y al tratamiento de la varianza, y la
reproducibilidad completa del análisis.

Entre las limitaciones, primero, la ENDES es autoinformada en varios
componentes, por lo que puede existir error de recuerdo en la edad
gestacional y el número de visitas. Segundo, la definición de parto
institucional combinó lugar del parto y personal que lo atendió y puede
diferir del indicador oficial del INEI; sin embargo, la definición
basada solo en el lugar produjo resultados prácticamente idénticos.
Tercero, algunas categorías de M13 y M15 no fueron clasificables y se
trataron como faltantes, aunque representaron menos del 1 \% de la
muestra. Cuarto, se detectaron estratos con una sola unidad primaria de
muestreo; las dos alternativas de estimación de varianza produjeron
conclusiones iguales. Quinto, la agregación de encuestas anuales no
constituye un panel. Por último, los modelos no incluyeron distancia a
establecimientos, barreras lingüísticas percibidas, aseguramiento ni
calidad de las consultas, por lo que los hallazgos deben interpretarse
como asociaciones observacionales y no como efectos de la lengua sobre
la atención.

\section{Implicaciones y
recomendaciones}\label{implicaciones-y-recomendaciones}

Los resultados justifican mantener sistemas de seguimiento de
indicadores maternos desagregados por lengua materna y residencia. Un
monitoreo útil debería presentar porcentajes e intervalos de confianza,
distinguir entre acceso y calidad, informar diferencias tanto relativas
como absolutas y documentar definiciones operativas comparables entre
años. Se recomienda incorporar el indicador de ocho o más controles y el
inicio en el primer trimestre como metas de seguimiento. En los
establecimientos, la información comprensible, la pertinencia cultural y
el trato respetuoso son dimensiones que las cifras de cobertura no
capturan. La norma peruana de atención del parto vertical con
pertinencia intercultural \citep{minsa2016} y los estudios en Puno
proporcionan un marco para orientar estas acciones, sin que esta
investigación haya evaluado su impacto. Futuras investigaciones podrían
incorporar distancia geográfica, aseguramiento y calidad del control
prenatal, así como métodos cualitativos que expliquen la persistencia de
la desigualdad residual en el parto institucional.

\section{Conclusiones}\label{conclusiones}

Entre las mujeres encuestadas en Puno durante 2014--2025, la atención
prenatal y el parto institucional mostraron diferencias según la lengua
materna. Las mujeres quechuahablantes y aimarahablantes presentaron
mayor residencia rural, menor educación superior y mayor pobreza que las
castellanohablantes. Tras el ajuste sociodemográfico, ambos grupos
indígenas mostraron menores probabilidades de parto institucional, y las
aimarahablantes también de inicio temprano del control prenatal y de
seis o más controles; la diferencia absoluta más amplia fue la del
inicio temprano. Estas asociaciones no disminuyeron significativamente
en el periodo reciente y resultaron robustas a definiciones y métodos
alternativos, lo que respalda su uso para orientar el monitoreo de la
equidad en los servicios maternos de Puno.

\section{Disponibilidad de datos y
código}\label{disponibilidad-de-datos-y-cuxf3digo}

Los microdatos originales están disponibles en el portal de microdatos
del INEI \citep{inei2025}. El código, los datos procesados y este
documento se encuentran en el repositorio del proyecto:
\url{https://github.com/USUARIO/ARTICULO_PRENATAL_ENDES_2014_2025}. Para
reproducir el informe basta abrir la carpeta del proyecto en Positron y
ejecutar \texttt{quarto render informe.qmd}. Este documento se generó el
23/11/2026 con R 4.3.3, Quarto 1.11.9 y los paquetes \texttt{survey}
4.2.1, \texttt{dplyr} 1.1.4, \texttt{ggplot2} 3.4.4, \texttt{knitr}
1.45.


\bibliography{referencias.bib}



\end{document}
